% Generated from config/requirements/section3_feasibility_v1.json.
\begin{table}[H]\centering\small\setlength{\tabcolsep}{3pt}
\caption{Mission Feasibility Requirements: System Integration}\label{tab:fr-sys}
\begin{tabular}{@{}
>{\RaggedRight\arraybackslash}p{\dimexpr(\linewidth-4\tabcolsep)*44/100\relax}
>{\RaggedRight\arraybackslash}p{\dimexpr(\linewidth-4\tabcolsep)*32/100\relax}
>{\RaggedRight\arraybackslash}p{\dimexpr(\linewidth-4\tabcolsep)*24/100\relax}
@{}}
\toprule Requirement & Simulation / Analysis Result & Feasibility Assessment\\\midrule
\phantomsection\label{req:ANT-SYS-FR-001}\mbox{ANT-SYS-}\newline\mbox{FR-001}\par The system feasibility assessment shall identify a common mass, epoch, orbit/RAAN, duty-cycle, response and resource baseline for each combined margin and disclose incompatible evidence rather than pool it. & Studies retain distinct epoch, mass and detector-response assumptions; no combined margin is inferred. & CONDITIONALLY\newline SUPPORTED\\
\addlinespace[3pt]
\bottomrule
\end{tabular}
\end{table}
\paragraph{Feasibility assessment.} The evidence supports scenario-level operations and resource feasibility. Each study retains its own configuration; the integrated mission conclusion in Section 5 distinguishes demonstrated results, conditional applications and unassessed constraints.
